\section{Leakage-Safe Splits, Protocols and the Shortcut Study}\label{sec:leakage}

\subsection{Leakage-Safe Splits}\label{sec:leakage-splits}

A tampered image shares most of its pixels with its host image. If the host is in the training set and the tampered
version in the test set, the evaluation is optimistic. We therefore split by \emph{groups}, built with union-find
over four kinds of link: the host id; exact pixel duplicates (196 merges); verified near duplicates (17 merges from
694 pairs); and the second-named id, for the 23 images where that image supplies most of the pixels (16 merges).
The result is 7,262 split groups of at most 69 images (near-duplicate criteria in Appendix~\ref{app:leakage}).

Stratified group $k$-fold assignment (20 bins, seed 42) with strata of forgery type $\times$ file format gives
train 8,828 / val 1,894 / test 1,892 images (70.0/15.0/15.0\%). The train+val (``dev'') set is further
divided into five grouped CV folds. No group crosses a split or a fold; this is asserted before the split is
written and checked again by the test suite. The split is frozen by SHA-256 hashes; it was re-frozen once, before
any model had seen the test set, after the forgery-type labels were corrected and the donor-dominant links added.
The test set stays unused until the final evaluation (Section~\ref{sec:test}).

One limitation remains. For 830 tampered images (828 splicing), the donor image is in a different split. The shared
content is the pasted region: a median of 5.5\% of the image, and 34.5\% at the 90th percentile. We report this
rather than prevent it: merging every donor would create a single group of 4,027 images (31.9\% of the dataset),
which makes a stratified 70/15/15 split impossible.

\subsection{Evaluation Protocols}\label{sec:leakage-protocols}

\emph{A (as released):} images exactly as distributed. \emph{C (JPEG-only):} A restricted to JPEG files of both
classes, an intuitive way to remove the format cue. \emph{B (re-encoded, Q=85):} every image is decoded and re-encoded
once to JPEG at quality 85 with 4:2:0 chroma subsampling, so all images share one container format and one
quantisation table. \emph{R (resampled + re-encoded, ``strict''):} every image is downscaled by 0.75 with area
interpolation, then re-encoded as in B; masks are resized with nearest-neighbour interpolation.

\subsection{Shortcut Study}\label{sec:leakage-shortcuts}

We train a random forest~\cite{breiman2001rf} (300 trees, balanced class weights) on feature families computed
\emph{globally over the whole image}, none of which localises anything: \emph{file metadata} (format flags,
quantisation-table statistics); \emph{image dimensions}; \emph{naive global ELA} (statistics of one ELA residual at
Q=90); and \emph{compression history} (global double-quantisation histograms of low-frequency DCT coefficients and
a global JPEG-ghost curve). Table~\ref{tab:leakage} and Fig.~\ref{fig:leakage} give balanced accuracy on the five
grouped dev folds (chance is 0.5); the full feature lists and ROC-AUCs are in Appendix~\ref{app:leakage}.

\begin{table}[t]
\centering
\caption{Global shortcut features alone: balanced accuracy (5-fold grouped CV on dev) per protocol}\label{tab:leakage}
\footnotesize
\begin{tabular}{lcccc}
\toprule
Features & A & C & B & R \\
\midrule
File metadata & \textbf{0.990} & \textbf{0.991} & 0.508 & 0.503 \\
Image dimensions & 0.603 & 0.597 & 0.603 & 0.604 \\
Naive global ELA & 0.696 & 0.678 & 0.564 & 0.539 \\
Compression history & 0.928 & 0.955 & \textbf{0.789} & 0.608 \\
All shortcuts & 0.989 & 0.986 & \textbf{0.797} & \textbf{0.649} \\
\bottomrule
\end{tabular}
\end{table}

\begin{figure*}[t]
\centering
\includegraphics[width=0.82\textwidth]{fig_leakage.png}
\caption{Balanced accuracy (mean and standard deviation over the five grouped dev folds) of a random forest trained
only on global shortcut features, under protocols A, C, B and R. The dashed line marks chance (0.5).}
\label{fig:leakage}
\end{figure*}

\textbf{The benchmark as released can be solved almost perfectly without looking at image content.} File metadata
alone reaches 99.0\% balanced accuracy (AUC 0.999). \textbf{Restricting to JPEG files does not fix this}: the
quantisation tables still separate the classes (99.1\%). \textbf{Re-encoding (B) removes the container/metadata
shortcut (0.508) but not the compression history.} A global double-quantisation and ghost-curve classifier still
reaches 0.789 under B. Each image carries its earlier compression history: authentic images were last saved as
JPEGs with standard libjpeg tables (89\% of them), while tampered images were last saved as editing-software JPEGs
or as uncompressed TIFFs, and one extra re-encode does not erase that history. The effect is strongest when both
classes started as JPEG: 0.919 on JPEG sources alone, against 0.725 for authentic vs.\ TIFF-tampered.
\textbf{Resampling before re-encoding (R) erases most of it.} Compression history falls to 0.608, and all shortcuts
together reach 0.649 $\pm$ 0.011; most of what remains is the image-size shortcut. \textbf{Image size is a residual
shortcut (0.60) that no protocol removes}, so the forensic pipeline must never use raw image dimensions as
features. Finally, \textbf{naive global ELA is weak once the container is controlled} (0.56 under B). Even under A
it reaches only 0.70, and it gains nothing from TIFF files: it scores 0.68 under C, where no TIFFs remain. Its
signal is compression history, not tampering.

% Module table (Section V) placed here so that the full-width float appears next to Section V.
\begin{table*}[t]
\centering
\caption{The eight forensic modules, the DIP technique family each one uses, and the evidence it computes}\label{tab:modules}
\scriptsize
\begin{tabular}{@{}c >{\raggedright\arraybackslash}p{2.3cm} >{\raggedright\arraybackslash}p{1.9cm} >{\raggedright\arraybackslash}p{8.6cm} >{\raggedright\arraybackslash}p{2.4cm}@{}}
\toprule
\# & Module & DIP category & Evidence & Main parameters \\
\midrule
1 & Error Level Analysis & Point processing &
$|I - \mathrm{JPEG}_q(I)|$ averaged per $8\times8$ block, divided by (block texture $+\,c$) so busy regions are not
mistaken for tampering & $q$, block, $c$ \\
2 & Patch-histogram $\chi^2$ & Histogram processing &
ELA residual contrast-stretched and histogram-equalised; $\chi^2$ distance of each patch histogram to the image's
reference histogram (the mean over textured patches) & $q$, patch, bins \\
3 & Noise level & Spatial filtering &
Immerk{\ae}r high-pass kernel (or a $3\times3$ Laplacian); per-block robust $\sigma = 1.4826\cdot\mathrm{median}|r|$
over non-edge pixels (gradient $\leq$ a global quantile), since edges leak content into the residual; flat and
clipped blocks excluded via the common validity mask & kernel, block, edge quantile \\
4 & JPEG ghosts & Point processing &
Per-block squared re-compression error at Q = 50 \ldots 95; each block's curve is min-max normalised; evidence =
L1 distance from the image's median curve & block \\
5 & DCT double quantisation & Frequency domain &
Per $8\times8$ DCT frequency: last quantisation step (known or estimated), test of the coefficient histogram for
periodicity, and per-coefficient posterior of \emph{not} following the periodic pattern; block log-odds averaged
over frequencies & number of frequencies, period threshold \\
6 & Edge sharpness & Edge detection &
Canny edges; at each edge pixel, sharpness = gradient / local contrast ($5\times5$ max $-$ min); block median; both
over-sharp (hard cut-out) and over-soft (feathered) outliers count & Canny thresholds, block \\
7 & Copy-move: keypoints & Corner/keypoint detection &
SIFT (or Harris) keypoints with SIFT descriptors; g2NN matching of the image against itself and against its mirror
image; pairs $<20$\,px apart dropped; RANSAC affine fit, separately for direct and mirrored matches & detector,
ratio, contrast \\
8 & Copy-move: blocks & Frequency domain &
Overlapping $16\times16$ blocks described by 9 low-frequency DCT coefficients; the blocks of the image and of its
mirror are sorted lexicographically together; near-identical neighbours are paired and their shift vectors voted on
& tolerance, min.\ block std, min.\ votes \\
\bottomrule
\end{tabular}
\end{table*}

\subsection{Consequences for the Rest of the Project}\label{sec:leakage-consequences}

A forensic feature set shows real tampering evidence only if it clearly beats the shortcuts under the same
protocol: 0.797 $\pm$ 0.007 balanced accuracy (AUC 0.876) under B and 0.649 (AUC 0.720) under R. These are lower
bounds; a stronger global probe might score higher. We therefore report classification under both B and R,
measure the \emph{added value} of the forensic features over the shortcuts, and choose the main protocol with a
rule fixed before the classification results were seen (Section~\ref{sec:classification}). Pixel-level
localisation cannot be gamed by global shortcuts, so we treat it as the most trustworthy result, and the features
favour \emph{within-image inconsistency} over absolute global statistics (Section~\ref{sec:method}). For B we keep
Q=85, which suppresses the shortcuts about as well as Q=75 while discarding less image detail
(Appendix~\ref{app:leakage}).
